\documentclass[12pt]{article}

\usepackage{amsmath, amssymb}
\usepackage{geometry}
\usepackage{setspace}
\usepackage{enumitem}

\geometry{margin=1in}
\setstretch{1.2}

\begin{document}

\begin{center}
\Large \textbf{CSE400 – Fundamentals of Probability in Computing} \\
\large \textbf{Lecture 6: Discrete Random Variables, Expectation and Problem Solving} \\
\normalsize (Exam-ready lecture scribe)
\end{center}

\vspace{1em}

Source: Lecture slides “CSE400 – Fundamentals of Probability in Computing, Lecture 6: Discrete RVs, Expectation and Problem Solving” by Dhaval Patel, PhD.

\section*{1. Random Variables (RVs)}

\subsection*{1.1 Motivation and Concept}

A \textbf{random variable (RV)} is a function that assigns a numerical value to each outcome in a sample space of a random experiment.

\begin{itemize}
\item The \textbf{distribution} of a random variable can be visualized using a \textbf{bar diagram}.
\item The \textbf{x-axis} represents the values the random variable can take.
\item The \textbf{height of the bar} at value $a_i$ represents the probability
\[
\Pr[X = a_i].
\]
\item Each probability is obtained by computing the probability of the \textbf{corresponding event} in the sample space.
\end{itemize}

\section*{2. Types of Random Variables}

Random variables are broadly classified into:

\subsection*{2.1 Discrete Random Variables}

A random variable is \textbf{discrete} if it can take \textbf{at most a countable number of values}.

\textbf{Key properties (as per slides):}
\begin{itemize}
\item Countable support
\item Probability Mass Function (PMF)
\item Probabilities assigned to \textbf{single values}
\item Each possible value has \textbf{strictly positive probability}
\end{itemize}

\subsection*{2.2 Continuous Random Variables}

A random variable is \textbf{continuous} if it takes values from an \textbf{uncountable set}.

\textbf{Key properties (as per slides):}
\begin{itemize}
\item Uncountable support
\item Probability Density Function (PDF)
\item Probabilities assigned to \textbf{intervals}
\item Each individual value has \textbf{zero probability}
\end{itemize}

(Note: This lecture focuses primarily on discrete random variables.)

\section*{3. Example of a Discrete Random Variable}

\subsection*{Example 1: Tossing Three Fair Coins}

\textbf{Experiment:} Toss 3 fair coins. \\
Let
\[
Y = \text{number of heads}.
\]

\textbf{Possible values of $Y$:}
\[
\{0, 1, 2, 3\}
\]

\textbf{Probabilities:}
\begin{itemize}
\item $P(Y = 0) = P(T,T,T) = \frac{1}{8}$
\item $P(Y = 1) = P(T,T,H) + P(T,H,T) + P(H,T,T) = \frac{3}{8}$
\item $P(Y = 2) = P(T,H,H) + P(H,T,H) + P(H,H,T) = \frac{3}{8}$
\item $P(Y = 3) = P(H,H,H) = \frac{1}{8}$
\end{itemize}

Thus,
\[
P(Y=0)=\frac{1}{8},\quad
P(Y=1)=\frac{3}{8},\quad
P(Y=2)=\frac{3}{8},\quad
P(Y=3)=\frac{1}{8}.
\]

Since $Y$ must take one of these values,
\[
\sum_{i=0}^{3} P(Y=i) = 1.
\]

This confirms that the probabilities form a valid distribution.

\section*{4. Probability Mass Function (PMF)}

\subsection*{4.1 Definition}

A random variable that can take \textbf{at most a countable number of possible values} is said to be \textbf{discrete}.

Let $X$ be a discrete random variable with range (possible values)
\[
R_X = \{x_1, x_2, x_3, \dots\}
\]
(finite or countably infinite).

The function
\[
P_X(x_k) = P(X = x_k), \quad k = 1,2,3,\dots
\]
is called the \textbf{Probability Mass Function (PMF)} of $X$.

\subsection*{4.2 Fundamental Property of PMF}

Since $X$ must take \textbf{one of the values} in its range,
\[
\sum_{k=1}^{\infty} P_X(x_k) = 1.
\]

This condition is \textbf{necessary and sufficient} for a valid PMF.

\section*{5. PMF – Solved Example}

\subsection*{Example: Given PMF with Unknown Constant}

The probability mass function of a random variable $X$ is given by
\[
p(i) = c \frac{\lambda^i}{i!}, \quad i = 0,1,2,\dots
\]
where $\lambda > 0$.

\textbf{Tasks:}
\begin{enumerate}
\item Find $P(X=0)$
\item Find $P(X>2)$
\end{enumerate}

\subsection*{Step 1: Find the Constant $c$}

Using the PMF property:
\[
\sum_{i=0}^{\infty} p(i) = 1
\]

Substitute:
\[
\sum_{i=0}^{\infty} c \frac{\lambda^i}{i!} = 1
\]

Factor out $c$:
\[
c \sum_{i=0}^{\infty} \frac{\lambda^i}{i!} = 1
\]

Using the known series:
\[
e^{\lambda} = \sum_{i=0}^{\infty} \frac{\lambda^i}{i!}
\]

So,
\[
c e^{\lambda} = 1
\quad \Rightarrow \quad
c = e^{-\lambda}
\]

\subsection*{Step 2: Find $P(X=0)$}

\[
P(X=0) = c \frac{\lambda^0}{0!}
= e^{-\lambda}.
\]

\subsection*{Step 3: Find $P(X>2)$}

\[
P(X>2) = 1 - P(X \le 2)
\]

\[
= 1 - [P(X=0) + P(X=1) + P(X=2)]
\]

Substitute values:
\[
= 1 - \left[
e^{-\lambda}
+ \lambda e^{-\lambda}
+ \frac{\lambda^2}{2} e^{-\lambda}
\right].
\]

\section*{6. Bayes’ Theorem (Recap)}

Using conditional probability:
\[
\Pr(A \mid B) = \frac{\Pr(B \mid A)\Pr(A)}{\Pr(B)}.
\]

For events $B_1, B_2, \dots, B_n$:
\[
\Pr(B_i \mid A) =
\frac{\Pr(A \mid B_i)\Pr(B_i)}
{\sum_{j=1}^{n} \Pr(A \mid B_j)\Pr(B_j)}.
\]

This is known as \textbf{Bayes’ Formula (Proposition 3.1)}.

\textbf{Terminology}
\begin{itemize}
\item $\Pr(B_i)$: \textbf{Prior probability}
\item $\Pr(B_i \mid A)$: \textbf{Posterior probability}
\end{itemize}

\section*{7. Bayes’ Theorem – Example 1}

\subsection*{Auditorium with 30 Rows}

\begin{itemize}
\item Auditorium has \textbf{30 rows}
\item Row 1 has 11 seats, Row 2 has 12 seats, …, Row 30 has 40 seats
\item A \textbf{row is selected uniformly} at random
\item A \textbf{seat within that row} is selected uniformly
\end{itemize}

\textbf{Tasks:}
\begin{enumerate}
\item Compute $P(\text{Seat 15} \mid \text{Row 20})$
\item Compute $P(\text{Row 20} \mid \text{Seat 15})$
\end{enumerate}

\subsection*{Solution Outline (as per slides)}

Since Row 20 has 30 seats:
\[
P(\text{Seat 15} \mid \text{Row 20}) = \frac{1}{30}.
\]

Total probability of selecting Seat 15:
\[
P(S_{15}) =
\sum_{k=15}^{30}
\frac{1}{30} \cdot \frac{1}{(k+10)}.
\]

Apply Bayes’ formula:
\[
P(\text{Row 20} \mid \text{Seat 15})
=
\frac{P(\text{Seat 15} \mid \text{Row 20}) P(\text{Row 20})}
     {P(\text{Seat 15})}.
\]

\section*{8. Bayes’ Theorem – Example 2}

\subsection*{Communication System (Receiver)}

\begin{itemize}
\item Binary data $\{0,1\}$ is transmitted
\item Receiver may make errors:
\begin{itemize}
\item 0 may be detected as 1
\item 1 may be detected as 0
\end{itemize}
\item System is described by given \textbf{conditional probabilities}
\end{itemize}

(The example applies Bayes’ theorem directly to compute posterior probabilities based on observed outputs.)

\section*{Final Exam Notes}

\begin{itemize}
\item \textbf{Discrete RV $\rightarrow$ PMF}
\item \textbf{Sum of PMF = 1}
\item Always verify normalization before computing probabilities
\item Use \textbf{complement rule} for tail probabilities
\item Bayes’ theorem = conditional probability + total probability law
\item Clearly identify \textbf{prior}, \textbf{likelihood}, and \textbf{posterior}
\end{itemize}

This scribe is \textbf{complete, faithful, and revision-safe}, and can be used as a \textbf{standalone reference} for this lecture.

\end{document}

